\chapter*{Abstract (EN)}
\addcontentsline{toc}{chapter}{Abstract (EN)}

\textbf{Title:} Development of an AI-Powered Telegram Bot for Personalized Gift Recommendations Using Prompt Engineering Techniques.

Choosing a gift is a common but time-consuming task in daily life. Online marketplaces show many products but no personalized advice. Modern Large Language Models (LLMs) can reason about a receiver profile and suggest gift ideas in natural language. However, the quality of LLM answers depends a lot on the prompt design. This problem is studied by the new research area called prompt engineering.

The aim of the thesis is to develop a Telegram bot that uses an LLM (Grok by xAI) for personalized gift recommendations, and to compare four prompt engineering strategies in this task: Naive, Constrained, Chain-of-Thought (CoT), and Persona-based.

The bot was developed in Python using the aiogram framework and SQLite as the database. The user passes through a short questionnaire with six structured questions, then receives three personalized gift ideas. The system was deployed on the PythonAnywhere cloud platform.

A controlled experiment was carried out. Thirty synthetic user personas were prepared. All four prompt strategies were run on all personas, giving 120 test cases. Each output was scored on five metrics: relevance, creativity, specificity, hallucination rate, and response time. The bot was also tested by eight real users.

The results confirm the hypothesis that advanced strategies produce significantly better recommendations than naive prompting. The best overall strategy was Chain-of-Thought (relevance 4.43/5, hallucination rate 8.9\%). The most creative strategy was Persona-based (creativity 4.53/5). All four strategies produced answers in less than 10 seconds. 75\% of real users said they would use the bot again.

The thesis also proposes future improvements, including real marketplace integration, multimodal input, multilingual support, and a hybrid prompt pipeline that combines Persona creativity with Constrained filtering.

The work confirms that prompt engineering is a real and measurable engineering activity. The methodology proposed in this thesis is domain-agnostic and can be applied to other personalized recommendation tasks.

\textbf{Volume:} 55 pages, 14 tables, 7 figures, 22 references, 3 appendices.
